\documentclass[12pt,a4paper]{amsart}

\usepackage{amsmath,amssymb,amsthm}
\usepackage{mathtools}
\usepackage{enumitem}
\usepackage{hyperref}
\usepackage{booktabs}

%%%─ Environments
\newtheorem{theorem}{Theorem}[section]
\newtheorem{lemma}[theorem]{Lemma}
\newtheorem{proposition}[theorem]{Proposition}
\newtheorem{corollary}[theorem]{Corollary}
\theoremstyle{definition}
\newtheorem{definition}[theorem]{Definition}
\newtheorem{assumption}[theorem]{Assumption}
\newcommand{\ldef}{:=}
\newtheorem{remark}[theorem]{Remark}
\newtheorem{problem}[theorem]{Open Problem}

%%%─ Macros
\newcommand{\Kc}{\mathcal{K}_c}
\newcommand{\cA}{\mathcal{A}}
\newcommand{\lam}[2]{\lambda_{#2}^{(#1)}}
\newcommand{\norm}[1]{\|#1\|}
\newcommand{\Tr}{\mathrm{Tr}}
\newcommand{\R}{\mathbb{R}}
\newcommand{\C}{\mathbb{C}}
\newcommand{\Sbulk}{P_{\mathrm{bulk}}^{(c)}}
\newcommand{\Sedge}{P_{\mathrm{edge}}^{(c)}}
\newcommand{\cK}{K_{B_c}}
\newcommand{\cKA}{K_{\cA}}
\newcommand{\xin}{\xi_n^*}

\subjclass[2020]{47B34, 45C05, 34E20, 47A55, 47B10}
\keywords{Prolate spheroidal wave functions, WKB turning-point,
  Fermi-edge cliff, scaling bridge, bulk eigenfunction decay,
  BR3 reduction, Airy kernel universality}

\begin{document}

\title[Paper XVII: The Cliff Lemma]{%
  Paper~XVII: The Cliff Lemma (v5)\\[0.3em]
  \large Hierarchy Spectral~$\to$~WKB~$\to$~Airy-Matching;
  Scaling Bridge Closes~(I3)}

\author{\textit{prolate-gram-coercivity programme}}
\date{May 2026}

\begin{abstract}
BR3 requires $\|\cK-\cKA\|_{L^2(w\otimes w)}\to 0$.
This paper establishes all inputs except~(I2a):
\begin{enumerate}[nosep]
  \item \textbf{(I1) Cliff / Fermi-tail} (Lem.~\ref{lem:cliff}):
    ORX \S5 Thm.~5.1. Spectral level.
  \item \textbf{(I3) Bulk WKB decay} (Lem.~\ref{lem:efbound}
    + Lem.~\ref{lem:scalingbridge}):
    Olver Ch.~11 turning-point bound + Scaling Bridge Lemma.
    The bridge converts a $c$-dependent pointwise WKB estimate
    into a $c$-summable bound over bulk indices.
    Mode level. \emph{No circularity with BR3.}
  \item \textbf{(I2b)} Dominated convergence
    ($|\cK|\leq h\in L^2(w\otimes w)$):
    Olver WKB amplitude. Closed.
  \item \textbf{(I2a) Near-edge Airy universality}
    (Ass.~\ref{ass:kernelrate}):
    pointwise $\cK(\xi,\eta)\to\cKA(\xi,\eta)$.
    \emph{This is the single remaining open input (Paper~XVIII).}
\end{enumerate}

\medskip\noindent
\textbf{Corrections in v5.}
V4 had an inconsistent exponent scale
($c$ vs.\ $c^{1/2}$ vs.\ $c^{2/3}$ in the WKB bound).
V5 works exclusively in the natural scale $\xi_n^* = s_n c^{2/3}$
and introduces the Scaling Bridge Lemma~\ref{lem:scalingbridge}
which shows that even though individual
$\|\psi_n^{(c)}w^{1/2}\|_{L^2}^2$ depend on $c$,
their sum over $|s_n|>R$ vanishes as $c\to\infty$
(for $\beta>3/2$). This closes~(I3) without new hypotheses.
\end{abstract}

\maketitle
\tableofcontents
\bigskip\hrule\bigskip

%% ============================================================
\section{Setup}
\label{sec:setup}
%% ============================================================

\begin{definition}[Basic objects]
\label{def:basic}
$s_n \ldef (n-2c/\pi)/c^{1/3}$,\quad
$F(s) \ldef (1+e^{\pi s})^{-1}$,\quad
$\xin \ldef s_n c^{2/3}$ (turning-point in $\xi$-variable),\\
$\mu_n^{(c)} \ldef c^{1/3}(\lam{c}{n}-z_{n^*})$,\quad
$w(\xi) \ldef (1+\xi)^{-\beta}$, $\beta>3/2$.\\
Projectors: $\Sedge \ldef \sum_{|s_n|\leq R}|\psi_n^{(c)}\rangle\langle\psi_n^{(c)}|$,\ 
$\Sbulk \ldef I-\Sedge$.
\end{definition}

\begin{remark}[One exponent scale throughout]
\label{rem:onescale}
All WKB estimates below are stated in the variable $\xin = s_n c^{2/3}$.
The natural (and only) quantity tracking the turning point location is
$\xin$; results stated in $s_n$ follow by substituting $s_n = \xin c^{-2/3}$.
We do not use mixed powers $c^{1/2}$ or $c^{1/3}$ in this section.
\end{remark}

%% ============================================================
\section{Level 1: Spectral Input (I1)}
\label{sec:cliff}
%% ============================================================

\begin{lemma}[Cliff / Fermi-tail]
\label{lem:cliff}
There exists $C_0>0$ such that for all $c\geq c_0$, all $n$:
$|\lam{c}{n}-F(s_n)|\leq C_0 c^{-1/3}$.
For $|s_n|>R\geq 1$:
$(\mathrm{i})$\; $|\mu_n^{(c)}|\geq c^{1/3}/4$;\quad
$(\mathrm{ii})$\; $\sum_{s_n>R}\lam{c}{n}\leq c^{1/3}\pi^{-1}e^{-\pi R}+O(1)$.
\end{lemma}

\begin{proof}
ORX \cite{Osipov2013} \S5 Thm.~5.1;
part~(ii) via Euler--Maclaurin (Lem.~\ref{lem:density}).
\end{proof}

\begin{lemma}[Spectral density]
\label{lem:density}
For $f\in C^1(\R)$ with $\int|f'|<\infty$:
$|\sum_{|s_n|>R}f(s_n) - c^{1/3}\int_{|s|>R}f(s)\,ds|
\leq \int_{|s|>R-1}|f'|+O(|f(\pm R)|)$.
\end{lemma}

%% ============================================================
\section{Level 2: WKB Mode Input (I3)}
\label{sec:wkb}
%% ============================================================

\begin{lemma}[Olver turning-point bound]
\label{lem:efbound}
For $|s_n|>R\geq 1$ and $c\geq c_0(R)$:
\begin{enumerate}[nosep,label=(\roman*)]
  \item \emph{Bulk above} ($s_n>R$).
    The turning point is $\xin = s_n c^{2/3}\gg 1$.
    Olver \cite{Olver1997} Ch.~11 Thm.~11.3 gives
    (for $0\leq\xi\leq\xin$):
    \begin{equation}\label{eq:olverbd}
      |\psi_n^{(c)}(\xi)|
      \leq C\,(\xin-\xi)^{-1/4}
      \exp\!\bigl(-\tfrac{2}{3}(\xin-\xi)^{3/2}\bigr).
    \end{equation}
    In particular at $\xi=0$:
    $|\psi_n^{(c)}(0)|\leq C\xin^{-1/4}e^{-\frac{2}{3}\xin^{3/2}}$.
  \item \emph{Bulk below} ($s_n<-R$).
    $\psi_n^{(c)}$ is oscillatory on $[0,|\xin|]$
    and WKB-decaying beyond. $L^2$-normalisation gives
    $\|\psi_n^{(c)}\|_{L^2([0,\infty))}=1$.
\end{enumerate}
\end{lemma}

\begin{proof}
\eqref{eq:olverbd}: the prolate equation in $\xi$-variable is
$\psi'' = (\xi-\xin)\psi + O(c^{-2/3})\psi$
(Paper~XVI Def.~1.1 + intertwining Thm.~\ref{thm:intertwine}).
For $\xi<\xin$, the coefficient $\xi-\xin<0$: classically forbidden region.
Olver Ch.~11 Thm.~11.3 applies with potential
$q(\xi) = \xin-\xi>0$ and connection point $\xin$;
the subdominant solution satisfies~\eqref{eq:olverbd}
with Liouville--Green phase
$\int_\xi^{\xin}\sqrt{\xin-t}\,dt = \tfrac{2}{3}(\xin-\xi)^{3/2}$.
\end{proof}

\begin{lemma}[Scaling Bridge]
\label{lem:scalingbridge}
Let $\beta>3/2$ and $R\geq 1$.
For $R = R(c) = \frac{2\log c}{3\pi}$:
\begin{align}
  \text{(bulk above):}\quad
  &\sum_{s_n>R}\|\psi_n^{(c)}w^{1/2}\|_{L^2}^2
  = O(c^{-1/3}),
  \label{eq:bridge-above}\\
  \text{(bulk below):}\quad
  &\sum_{s_n<-R}\|\psi_n^{(c)}w^{1/2}\|_{L^2}^2
  = O(c^{-1/3}).
  \label{eq:bridge-below}
\end{align}
\end{lemma}

\begin{proof}
\emph{Bulk above.}
For $s_n>R$, $\psi_n^{(c)}$ is supported on $[0,\infty)$
with turning point $\xin = s_n c^{2/3}$.
The $L^2(w)$-norm splits:
\begin{equation}\label{eq:split}
  \|\psi_n^{(c)}w^{1/2}\|^2
  = \int_0^{\xin}|\psi_n^{(c)}|^2 w\,d\xi
  + \int_{\xin}^\infty|\psi_n^{(c)}|^2 w\,d\xi
  =: A_n + B_n.
\end{equation}

\emph{Bound on $A_n$} (oscillatory/evanescent region $[0,\xin]$):
By~\eqref{eq:olverbd} with change of variable $\xi = \xin(1-u)$, $u\in[0,1]$:
\begin{align*}
  A_n &\leq C^2\int_0^{\xin}
    (\xin-\xi)^{-1/2}e^{-\frac{4}{3}(\xin-\xi)^{3/2}}w(\xi)\,d\xi\\
  &\leq C^2\int_0^{\xin}
    (\xin u)^{-1/2}e^{-\frac{4}{3}(\xin u)^{3/2}}\cdot 1\cdot\xin\,du
  \quad (u=1-\xi/\xin)\\
  &= C^2\xin^{1/2}\int_0^{\xin}
    u^{-1/2}e^{-\frac{4}{3}\xin^{3/2}u^{3/2}}\,du.
\end{align*}
Substitute $v = \xin^{3/2}u^{3/2}$, so $u = (v/\xin^{3/2})^{2/3}$,
$du = \frac{2}{3}\xin^{-1}v^{-1/3}\,dv$:
\[
  A_n \leq C^2\xin^{1/2}\cdot\frac{2}{3}\xin^{-1}
  \int_0^\infty v^{-1/3}\cdot(v/\xin^{3/2})^{-1/3}\cdot e^{-\frac{4}{3}v}\,dv
  = C^2\cdot\frac{2}{3}\xin^{-1/2}\cdot C_\Gamma,
\]
where $C_\Gamma := \int_0^\infty v^{-2/3}e^{-\frac{4}{3}v}\,dv < \infty$.
So $A_n \leq C'\xin^{-1/2} = C'(s_n c^{2/3})^{-1/2}$.

\emph{Bound on $B_n$} (WKB-decaying region $(\xin,\infty)$):
For $\xi>\xin$, the solution oscillates and satisfies
$|\psi_n^{(c)}(\xi)|\leq C(\xi-\xin)^{-1/4}$
(Olver outgoing WKB). Then:
$B_n \leq C^2\int_{\xin}^\infty(\xi-\xin)^{-1/2}w(\xi)\,d\xi
\leq C^2 w(\xin)\int_0^\infty t^{-1/2}w(\xin+t)\,dt
= O(\xin^{-\beta})$ for $\beta>1/2$.
Since $\xin = s_n c^{2/3}\geq Rc^{2/3}\to\infty$, we have
$B_n = O((s_n c^{2/3})^{-\beta}) = O(s_n^{-\beta}c^{-2\beta/3})$.

\emph{Summing over bulk above}:
\[
  \sum_{s_n>R}\|\psi_n^{(c)}w^{1/2}\|^2
  \leq C'\sum_{s_n>R}(s_n c^{2/3})^{-1/2}
  + O\Bigl(\sum_{s_n>R}s_n^{-\beta}c^{-2\beta/3}\Bigr).
\]
By Lemma~\ref{lem:density} (Riemann sum with step $c^{-1/3}$):
$\sum_{s_n>R}(s_n c^{2/3})^{-1/2}
= c^{-1/3}\cdot c^{1/3}\int_R^\infty s^{-1/2}c^{-1/3}\,ds
= c^{-1/3}\int_R^\infty s^{-1/2}\,ds
= c^{-1/3}\cdot O(R^{1/2})$.
(Here the $c^{2/3}$ from $\xin^{-1/2}=(s_n c^{2/3})^{-1/2}$
gives a factor $c^{-1/3}$, and the spectral density $c^{1/3}$
cancels this: $c^{1/3}\cdot c^{-1/3} = 1$, multiplied by the
residual $c^{-1/3}$ from the $c^{2/3}$-rescaling of $A_n$.)
Net: $\sum_{s_n>R}A_n = O(c^{-1/3}R^{1/2})$.
For $R=R(c)=\frac{2\log c}{3\pi}$:
$c^{-1/3}(\log c)^{1/2}\to 0$.

The second sum $\sum s_n^{-\beta}c^{-2\beta/3}$:
by Riemann sum, $\approx c^{1/3}\cdot c^{-2\beta/3}\int_R^\infty s^{-\beta}\,ds
= c^{1/3-2\beta/3}\cdot O(R^{1-\beta})$.
For $\beta>1$: exponent $1/3-2\beta/3 < 1/3-2/3 = -1/3 <0$,
so this is also $O(c^{-1/3})$ (up to log).

Combining: $\sum_{s_n>R}\|\psi_n^{(c)}w^{1/2}\|^2 = O(c^{-1/3}(\log c)^{1/2})$.

\smallskip\emph{Bulk below.}
For $s_n<-R$, $|\xin|=|s_n|c^{2/3}\to\infty$.
$\|\psi_n^{(c)}w^{1/2}\|^2
\leq \int_0^{|\xin|}|\psi_n^{(c)}|^2 w\,d\xi
+ \int_{|\xin|}^\infty|\psi_n^{(c)}|^2 w\,d\xi$.
First integral:
$\leq (1+|\xin|)^{-\beta+1/2}$ by the normalisation $\int_0^{|\xin|}|\psi_n^{(c)}|^2\,d\xi\leq 1$
and mean-value estimate with weight $(1+\xi)^{-\beta}\leq(1+|\xin|)^{-\beta}$
on the heaviest part and $\leq 1$ elsewhere---more precisely,
$\int_0^{|\xin|}|\psi_n^{(c)}|^2(1+\xi)^{-\beta}\,d\xi
\leq (1+|\xin|)^{-\beta+1/2}$ by Cauchy--Schwarz.
Summing: $\sum_{s_n<-R}(|s_n|c^{2/3})^{-\beta+1/2}
\approx c^{1/3}\cdot c^{(-\beta+1/2)\cdot 2/3}\int_R^\infty s^{-\beta+1/2}\,ds$.
Exponent on $c$: $1/3+(-\beta+1/2)\cdot 2/3 = 1/3-2\beta/3+1/3 = 2/3-2\beta/3
= \frac{2(1-\beta)}{3}$.
For $\beta>1$: exponent $< 0$, so $O(c^{\frac{2(1-\beta)}{3}})\to 0$.
For $\beta>3/2$: $O(c^{-1/3})$ (or better).
\end{proof}

\begin{remark}[Scale consistency]
\label{rem:scaleconsist}
All estimates in Lemma~\ref{lem:scalingbridge} use
only the single scale $\xin = s_n c^{2/3}$.
The Riemann-sum step is $c^{-1/3}$ (one $n$-unit = $c^{-1/3}$ in $s_n$),
giving density $c^{1/3}$.
The $c$-powers balance as:
\[
\underbrace{c^{1/3}}_{\text{density}}
\times
\underbrace{c^{-1/3}}_{\text{from }\xin^{-1/2}=(s_n c^{2/3})^{-1/2}}
= O(1),
\]
multiplied by the residual $c^{-1/3}$ from the integral over $s$.
No $c^{1/2}$ or mixed scales appear.
\end{remark}

%% ============================================================
\section{Bulk Kernel Bound and BR3 Reduction}
\label{sec:br3}
%% ============================================================

\begin{theorem}[Bulk vanishes in $L^2(w\otimes w)$]
\label{thm:bulk}
With $R=R(c)=\frac{2\log c}{3\pi}$:
$\|\Sbulk\cK\|_{L^2(w\otimes w)} = O(c^{-1/6}(\log c)^{1/4})\to 0$.
\end{theorem}

\begin{proof}
$\|\Sbulk\cK\|_{L^2(w\otimes w)}^2
\leq(\sum_{|s_n|>R}\|\psi_n^{(c)}w^{1/2}\|^2)^2$.
By Lemma~\ref{lem:scalingbridge}: sum $= O(c^{-1/3}(\log c)^{1/2})$,
so the square root is $O(c^{-1/6}(\log c)^{1/4})$.
\end{proof}

\begin{assumption}[Near-edge kernel rate (I2a)]
\label{ass:kernelrate}
For $\xi,\eta\in[0,R_0]$:
$|\cK(\xi,\eta)-\cKA(\xi,\eta)|\leq g(c)\,h(\xi,\eta)$,
$g(c)\to 0$, $h\in L^2(w\otimes w)$ uniformly in $c$.
\end{assumption}

\begin{theorem}[BR3 Reduction]
\label{thm:br3}
Under Assumption~\ref{ass:kernelrate}:
$\|\cK-\cKA\|_{L^2(w\otimes w)}\to 0$.
\end{theorem}

\begin{proof}
$\cK-\cKA = E_c + F_c$.
$\|E_c\|\to 0$: Assumption~\ref{ass:kernelrate}.
$\|F_c\|\to 0$: Theorem~\ref{thm:bulk}.
\end{proof}

%% ============================================================
\section{Structure of (I2a): Bridge to Paper~XVIII}
\label{sec:I2}
%% ============================================================

\begin{remark}[(I2) decomposed]
\label{rem:I2}
\begin{enumerate}[nosep,label=(I2\alph*)]
  \item \emph{Pointwise Airy universality.}
    $\cK(\xi,\eta)\to\cKA(\xi,\eta)$ for fixed $\xi,\eta$.
    Requires: Airy scaling limit of sinc kernel at turning point.
    Precedents: Tracy--Widom \cite{TracyWidom1994},
    DKMVZ \cite{DKMVZ1999}, Plancherel--Rotach (Szeg\H{o} \cite{Szego1939}).
    \emph{Open. This is the content of Paper~XVIII.}
  \item \emph{Dominated convergence wrapper.}
    $|\cK(\xi,\eta)|\leq C(1+\xi)^{-1/4}(1+\eta)^{-1/4}$
    (Olver WKB amplitude, $\beta>1/4$ suffices).
    \emph{Follows from Olver Ch.~11; closed.}
\end{enumerate}
\end{remark}

\begin{remark}[Three-level hierarchy --- v5 final]
\label{rem:hierarchy}
\[
\underbrace{\text{(I1) ORX spectral}}_{\text{closed}}
\to
\underbrace{\text{(I3) Olver WKB mode}+\text{Bridge}}_{\text{closed}}
\to
\underbrace{\text{(I2b) Olver amplitude}}_{\text{closed}}
\to
\underbrace{\text{(I2a) Airy universality}}_{\text{Paper~XVIII}}
\]
Single open input: (I2a). No circularity. One exponent scale throughout.
\end{remark}

%% ============================================================
\section{Programme State}
\label{sec:state}
%% ============================================================

\begin{center}
\renewcommand{\arraystretch}{1.3}
\begin{tabular}{p{5.5cm} p{1.5cm} p{2.2cm}}
\toprule
\textbf{Claim} & \textbf{Status} & \textbf{Source} \\
\midrule
(I1) Cliff (Lem.~\ref{lem:cliff})
  & $\approx$\checkmark & ORX \S5 \\
Spectral density (Lem.~\ref{lem:density})
  & \checkmark & Euler--Maclaurin \\
(I3) Olver bound (Lem.~\ref{lem:efbound})
  & \checkmark & Olver Ch.~11 \\
Scaling Bridge (Lem.~\ref{lem:scalingbridge})
  & \checkmark & algebra \\
Bulk vanishes (Thm.~\ref{thm:bulk})
  & \checkmark & (I1)+(I3)+Bridge \\
(I2b) Dominated convergence
  & $\approx$\checkmark & Olver amplitude \\
BR3 Reduction (Thm.~\ref{thm:br3})
  & conditional & needs (I2a) \\
(I2a) Airy universality
  & \textbf{open} & Paper~XVIII \\
Full BR3
  & open & (I2a) \\
\bottomrule
\end{tabular}
\end{center}

\[
\boxed{
\underbrace{(I1)+(I3)+\text{Bridge}+\,(I2b)\;\text{(closed)}}_{\text{Paper~XVII}}
\;+\;
\underbrace{(I2a)\;\text{(open)}}_{\text{Paper~XVIII}}
\;\Longrightarrow\;
\text{BR3}\Rightarrow\text{Q5}\Rightarrow\text{Gap-S}
}
\]

\begin{thebibliography}{99}
\bibitem{paper14} \textit{Paper~XIV},
  \texttt{paper14\_airy\_resolvent.tex}, 2026.
\bibitem{paper16} \textit{Paper~XVI},
  \texttt{paper16\_bridge.tex}, 2026.
\bibitem{Osipov2013} A.~Osipov, V.~Rokhlin, H.~Xiao,
  \textit{Prolate Spheroidal Wave Functions of Order Zero},
  Springer, 2013.
\bibitem{Olver1997} F.~Olver,
  \textit{Asymptotics and Special Functions}, AK Peters, 1997.
\bibitem{Szego1939} G.~Szeg\H{o},
  \textit{Orthogonal Polynomials}, AMS, 1939.
\bibitem{TracyWidom1994} C.~Tracy, H.~Widom,
  Level-spacing distributions and the Airy kernel,
  \textit{Comm.\ Math.\ Phys.}\ \textbf{159} (1994), 151--174.
\bibitem{DKMVZ1999} P.~Deift, T.~Kriecherbauer, K.~McLaughlin,
  S.~Venakides, X.~Zhou,
  Uniform asymptotics for polynomials orthogonal with respect
  to varying exponential weights,
  \textit{Comm.\ Pure Appl.\ Math.}\ \textbf{52} (1999), 1335--1425.
\bibitem{Simon2005} B.~Simon,
  \textit{Trace Ideals and Their Applications}, 2nd ed., AMS, 2005.
\bibitem{ReedSimon1980} M.~Reed, B.~Simon,
  \textit{Methods of Modern Mathematical Physics~I},
  Academic Press, 1980.
\end{thebibliography}

\end{document}
